% !TEX program = xelatex
\documentclass{misisarticle}

\grnti{28.23.15}
\asjc{1702}
\oecd{1.2}

\title{Поиск дублирующихся рекламных креативов на основе визуальных и текстовых представлений}

\misisauthor{И.~С.~Проскурякова}{кафедра инженерной кибернетики НИТУ «МИСиС»}{Москва, Россия}{proskuryakova0404@mail.ru}
\misisauthor{Э.~Р.~Тагиев}{научный руководитель, кафедра инженерной кибернетики НИТУ «МИСиС»}{Москва, Россия}{}

\setabstract{Рассматривается задача поиска похожих рекламных креативов по изображениям и текстам. Выделяются две бинарные задачи: визуальное сходство и совпадение рекламируемого продукта. Описаны формирование пар из пользовательских сессий, получение разметки Gemma и извлечение признаков моделями DINOv3, multilingual-e5-large и CLIP. Для объединения признаков рассматриваются взвешенная сумма, логистическая регрессия и CatBoost с дополнительным расстоянием pHash. Представлена методика оценки с разделением данных по сессиям и учётом ограничений автоматической разметки.}

\setkeywords{рекламные креативы, поиск дубликатов, мультимодальные представления, DINOv3, E5, CLIP, CatBoost}

\begin{document}
\maketitle

\section{Введение}

В рамках одной пользовательской сессии онлайн-платформа может показывать несколько рекламных креативов из разных групп и кампаний. Часть объявлений повторяет одно изображение с небольшими изменениями, использует общий рекламный шаблон или продвигает один продукт с различающимися изображениями и описаниями. Выявление таких случаев необходимо для анализа повторяемости рекламного контента и последующего управления разнообразием показов.

Точное совпадение файлов не охватывает все виды повторяющегося контента. Изменение размера изображения, редактирование текста или замена рекламного шаблона могут сохранять сходство объявлений с точки зрения пользователя. При этом визуальная похожесть и совпадение рекламируемого продукта не эквивалентны: похожие баннеры могут продвигать разные товары, а один товар может быть представлен визуально различными креативами. Поэтому в работе эти признаки рассматриваются как отдельные целевые переменные.

Цель работы состоит в разработке метода попарного сравнения рекламных креативов на основе визуальных и текстовых представлений. Рассматриваются отдельные оценки сходства изображений и текстов, их взвешенное объединение и обучаемая метамодель. Предлагаемый подход использует компактный набор признаков сходства вместо непосредственного обучения классификатора на исходных изображениях и текстах.

\section{Анализ работ}

\subsection{Текстовые представления}

Для сопоставления текстов могут применяться лексические признаки, включая TF--IDF, и расстояния между строками. Они удобны для поиска совпадающих слов и небольших редакционных изменений, однако не задают явного представления смысловой близости при перефразировании.

В работе Sentence-BERT~\cite{sbert} предложено получать векторные представления предложений, которые можно сравнивать с помощью косинусной близости. Это позволяет вычислять представление каждого текста отдельно и повторно использовать его при сравнении разных пар. Для рекламных описаний такой подход удобен, поскольку одни и те же креативы могут встречаться во многих сессиях.

Семейство Multilingual E5~\cite{e5} использует контрастивное предобучение на многоязычных парах текстов и последующее обучение на размеченных данных. В настоящей работе применяется multilingual-e5-large. Текстовое представление позволяет учитывать сведения о продукте, присутствующие в заголовке и описании, но его информативность зависит от полноты этих полей и наличия шаблонных фрагментов.

\subsection{Визуальные представления}

Перцептивное хэширование преобразует изображение в компактный код для поиска одинаковых и близких изображений~\cite{hashing}. Расстояние между кодами удобно использовать как дополнительный признак почти точного совпадения. При этом устойчивость к преобразованиям зависит от алгоритма хэширования: совпадение продукта на разных фотографиях само по себе не гарантирует близость хэшей.

Альтернативой являются визуальные эмбеддинги нейронных сетей. В DINOv3~\cite{dinov3} предложены визуальные представления, полученные с помощью самоконтролируемого обучения. В рассматриваемом методе используется вариант на основе Vision Transformer с представлением специального токена CLS. Сходство таких векторов применяется для сравнения изображений, в том числе не являющихся точными копиями. Однако близость по композиции или рекламному шаблону не обязательно означает совпадение продукта.

В работе Jakhar и Borah~\cite{near_duplicates} перцептивное хэширование объединяется с сиамской сетью и Vision Transformer для поиска почти дублирующихся изображений. Эта работа служит примером совместного использования компактных признаков и нейросетевых представлений, но решает задачу сравнения изображений без анализа текстовых полей рекламного объявления.

\subsection{Мультимодальные подходы}

Модель CLIP~\cite{clip} содержит визуальный и текстовый энкодеры, обучаемые контрастивно на парах «изображение~--- текст». Их выходы согласуются в общем пространстве представлений. Энкодеры обучаются совместно, но каждый эмбеддинг вычисляется по одной модальности. В работе они дают две дополнительные оценки: сходства изображений и сходства текстов.

При раннем объединении текстовые и визуальные признаки соединяются до принятия решения, например конкатенацией эмбеддингов. При позднем объединении, или late fusion, сначала независимо вычисляются оценки по каждой модальности, а затем формируется общий результат. Последний вариант позволяет использовать специализированные энкодеры и отдельно анализировать вклад каждого сигнала.

В качестве обучаемого объединения рассматриваются логистическая регрессия и градиентный бустинг CatBoost~\cite{catboost}. Первая задаёт линейную комбинацию входных признаков до применения логистической функции, второй позволяет моделировать нелинейные зависимости между оценками сходства. Качество способов объединения оценивается отдельно для каждой целевой переменной.

\section{Постановка задачи}

\subsection{Входные данные и целевые переменные}

Рекламный креатив представляется объектом $a_i=(I_i,T_i)$, где $I_i$~--- изображение, а $T_i$~--- текст, сформированный из заголовка, описания и юридического описания объявления. На вход метода поступает пара $(a_i,a_j)$, на выходе формируется оценка сходства относительно выбранного критерия.

Для каждой пары рассматриваются две независимые бинарные метки: $y_{ij}^{(v)}$ для визуального сходства и $y_{ij}^{(p)}$ для совпадения рекламируемого продукта или услуги. Для каждой задачи обучается отдельная метамодель и выбирается собственный порог. Решение записывается в виде
\begin{equation}
  \widehat y_{ij}^{(r)} =
  \mathbb{I}\!\left[s_r(a_i,a_j)\geq\tau_r\right],
  \quad r\in\{v,p\},
  \label{eq:decision}
\end{equation}
где $s_r$~--- оценка сходства, $\tau_r$~--- порог, а $\mathbb{I}$~--- индикатор выполнения условия. Единой агрегированной метки дубля в данной постановке не вводится.

\subsection{Формирование пар и разметка}

Исходные данные содержат пользовательские сессии с последовательностями показанных рекламных креативов. Из каждой сессии удаляются повторные показы одного и того же креатива. Для сессий, содержащих от двух до десяти различных креативов, строятся все неупорядоченные пары. Сессия длины $n$ формирует $n(n-1)/2$ пар.

В исходном исследовании набор для сессий длиной до пяти креативов содержит около 166~тыс. пар, а набор для сессий длиной от шести до десяти креативов~--- около 250~тыс. пар. Общий объём составляет более 416~тыс. пар до исключения неопределённых ответов по каждой целевой переменной. Это число относится к парам объявлений, а не к уникальным креативам или сессиям.

Разметка получена с помощью модели семейства Gemma. Для каждой пары задаются признаки визуального сходства, совпадения продукта и совпадения рекламодателя со значениями \texttt{yes}, \texttt{no} или \texttt{uncertain}. Первые два признака используются как целевые переменные. Значения \texttt{yes} и \texttt{no} преобразуются в 1 и 0; пары с ответом \texttt{uncertain} исключаются из обучения и оценки соответствующей задачи. Признак совпадения рекламодателя сохраняется в данных, но целевой переменной не является.

Автоматическая разметка может содержать систематические ошибки. Без независимой экспертной проверки полученные метрики следует интерпретировать как согласие с разметкой Gemma. Исключение неопределённых ответов также ограничивает оценку подмножеством пар, для которых модель смогла сформировать определённую метку.

\section{Описание методов}

\subsection{Визуальное и текстовое сходство}

Изображение обрабатывается предобученным энкодером DINOv3. Визуальным представлением служит выходной вектор токена CLS. Текстовые поля объединяются в одно описание и преобразуются энкодером multilingual-e5-large. После извлечения признаки нормализуются по евклидовой норме:
\begin{equation}
  z_i^{(m)}=\frac{f_m(x_i^{(m)})}
  {\left\|f_m(x_i^{(m)})\right\|_2},
  \quad m\in\{I,T\},
  \label{eq:normalization}
\end{equation}
где $x_i^{(I)}=I_i$, $x_i^{(T)}=T_i$, а $f_I$ и $f_T$~--- визуальный и текстовый энкодеры соответственно.

Для каждой модальности вычисляется косинусная близость, которая для нормализованных векторов равна скалярному произведению:
\begin{equation}
  s_m(a_i,a_j)=\left(z_i^{(m)}\right)^{\top}z_j^{(m)}.
  \label{eq:cosine}
\end{equation}
Оценки $s_I$ и $s_T$ используются как самостоятельные базовые методы и как признаки для объединения модальностей. Их значения не интерпретируются непосредственно как вероятности совпадения продукта или визуального сходства.

\subsection{Позднее объединение оценок}

В базовом мультимодальном подходе итоговая оценка задаётся взвешенной суммой:
\begin{equation}
  s_r^{\mathrm{fusion}}=
  w_{I,r}\,s_I+w_{T,r}\,s_T,
  \quad r\in\{v,p\}.
  \label{eq:fusion}
\end{equation}
Весовые коэффициенты и порог выбираются отдельно для визуального сходства и совпадения продукта на валидационной выборке. Такая схема сохраняет независимое вычисление эмбеддингов и не требует дополнительного обучения самих энкодеров.

\subsection{Дополнительные признаки CLIP и pHash}

Энкодеры CLIP формируют нормализованные представления изображения и текста. Между двумя изображениями вычисляется оценка $c_I$, а между двумя текстами~--- $c_T$ по правилу~\eqref{eq:cosine}. Эти оценки дополняют признаки DINOv3 и E5; межмодальные сравнения изображения одного объявления с текстом другого в описываемую схему не включаются.

Для выявления почти точных визуальных совпадений применяется перцептивный хэш pHash. Пусть $h(I)$~--- бинарный хэш изображения длины $L$. Тогда признак пары определяется расстоянием Хэмминга:
\begin{equation}
  d_H(I_i,I_j)=\sum_{k=1}^{L}
  \mathbb{I}\!\left[h(I_i)_k\neq h(I_j)_k\right].
  \label{eq:phash}
\end{equation}
Меньшее расстояние означает большее сходство хэшей. Этот признак имеет противоположное направление по отношению к косинусным оценкам и передаётся обучаемой метамодели как отдельная величина.

\subsection{Обучаемое объединение признаков}

Для пары креативов формируется вектор
\begin{equation}
  \boldsymbol\phi_{ij}=
  \left(s_I,\ s_T,\ c_I,\ c_T,\ d_H\right)^{\top}.
  \label{eq:features}
\end{equation}
В линейном варианте используются четыре оценки сходства $\boldsymbol\psi_{ij}=(s_I,s_T,c_I,c_T)^{\top}$. Вероятность положительного класса оценивается логистической регрессией:
\begin{equation}
  p_r=\sigma\!\left(b_r+
  \boldsymbol\beta_r^{\top}\boldsymbol\psi_{ij}\right),
  \qquad \sigma(t)=\frac{1}{1+e^{-t}}.
  \label{eq:logistic}
\end{equation}
Коэффициенты $\boldsymbol\beta_r$ и свободный член $b_r$ обучаются отдельно для каждой задачи. В нелинейном варианте CatBoost получает четыре оценки сходства или полный вектор~\eqref{eq:features} с дополнительным расстоянием pHash. Это позволяет учитывать сочетания сигналов, например высокое визуальное сходство при различающихся текстах. Предсказанная вероятность используется как оценка $s_r$ в правиле~\eqref{eq:decision}.

Для сопоставления способов объединения рассматриваются наборы из двух признаков DINOv3 и E5, из четырёх признаков с добавлением CLIP и из пяти признаков с добавлением pHash. Сравнение таких конфигураций позволяет оценивать вклад дополнительных источников информации на одной и той же выборке соответствующей задачи.

\subsection{Методика оценки качества}

Данные разделяются по идентификатору пользовательской сессии на обучающую, валидационную и тестовую части в соотношении 64/16/20. Пары из одной сессии относятся только к одной части. Разделение по сессиям само по себе не исключает появления одинаковых креативов в разных частях набора. Для оценки переноса на ранее не встречавшиеся креативы требуется отдельная проверка их пересечения или дополнительный протокол группировки данных.

Метамодели обучаются на обучающей части. Веса late fusion, параметры моделей, ранняя остановка и пороги классификации выбираются на валидационной части. Критерием выбора порога служит максимум F1-меры. Тестовая часть используется для итогового сравнения после фиксации конфигурации.

Качество ранжирования оценивается с помощью ROC-AUC и PR-AUC. Для порогового решения~\eqref{eq:decision} рассчитываются точность $P$, полнота $R$ и F1-мера:
\begin{equation}
  P=\frac{TP}{TP+FP},\qquad
  R=\frac{TP}{TP+FN},
\end{equation}
\begin{equation}
  F_1=\frac{2PR}{P+R}
      =\frac{2TP}{2TP+FP+FN},
\end{equation}
где $TP$, $FP$ и $FN$~--- числа истинно положительных, ложно положительных и ложно отрицательных решений. При сравнении методов необходимо сохранять одинаковый состав тестовых пар внутри каждой задачи и указывать долю положительного класса. Размеры выборок для двух задач могут различаться из-за исключения ответов \texttt{uncertain}.

Выбранная единица оценки~--- пара креативов из пользовательской сессии. Полученные по этому протоколу показатели не характеризуют автоматически полноту поиска кандидатов во всей базе объявлений или влияние метода на рекламные показатели.

\begin{thebibliography}{99}

\bibitem{sbert}
Reimers N., Gurevych I. Sentence-BERT: Sentence Embeddings using Siamese BERT-Networks. \textit{Proceedings of EMNLP-IJCNLP}, 2019, pp.~3982--3992. \url{https://aclanthology.org/D19-1410/}.

\bibitem{e5}
Wang L., Yang N., Huang X., Yang L., Majumder R., Wei F. Multilingual E5 Text Embeddings: A Technical Report. \textit{arXiv:2402.05672}, 2024. \url{https://arxiv.org/abs/2402.05672}.

\bibitem{hashing}
Biswas R., Blanco-Medina P. State of the Art: Image Hashing. \textit{arXiv:2108.11794}, 2021. \url{https://arxiv.org/abs/2108.11794}.

\bibitem{dinov3}
Sim\'eoni O., Vo H.~V., Seitzer M. et al. DINOv3. \textit{arXiv:2508.10104}, 2025. \url{https://arxiv.org/abs/2508.10104}.

\bibitem{near_duplicates}
Jakhar Y., Borah M.~D. Effective near-duplicate image detection using perceptual hashing and deep learning. \textit{Information Processing \& Management}, 2025, vol.~62, no.~4, article~104086. \url{https://doi.org/10.1016/j.ipm.2025.104086}.

\bibitem{clip}
Radford A., Kim J.~W., Hallacy C. et al. Learning Transferable Visual Models From Natural Language Supervision. \textit{Proceedings of ICML}, 2021, vol.~139, pp.~8748--8763. \url{https://proceedings.mlr.press/v139/radford21a.html}.

\bibitem{catboost}
Dorogush A.~V., Ershov V., Gulin A. CatBoost: gradient boosting with categorical features support. \textit{arXiv:1810.11363}, 2018. \url{https://arxiv.org/abs/1810.11363}.

\end{thebibliography}
\end{document}
